% !TeX program = pdflatex
\documentclass[11pt,a4paper]{amsart}

\newcommand{\notelanguage}{finnish}
% Shared preamble for course notes and exercise sets.

\usepackage[T1]{fontenc}
\usepackage{lmodern}
\providecommand{\notelanguage}{english}
\usepackage[finnish,english]{babel}
\AtBeginDocument{\expandafter\selectlanguage\expandafter{\notelanguage}}

\newcommand{\notesfinnishlanguage}{finnish}
\ifx\notelanguage\notesfinnishlanguage
\newcommand{\notestheoremname}{Lause}
\newcommand{\noteslemmaname}{Lemma}
\newcommand{\notespropositionname}{Propositio}
\newcommand{\notescorollaryname}{Korollaari}
\newcommand{\notesdefinitionname}{Määritelmä}
\newcommand{\notesexamplename}{Esimerkki}
\newcommand{\notesproblemname}{Tehtävä}
\newcommand{\notesexercisename}{Tehtävä}
\newcommand{\notesremarkname}{Huomautus}
\newcommand{\notessolutionname}{Ratkaisu}
\else
\newcommand{\notestheoremname}{Theorem}
\newcommand{\noteslemmaname}{Lemma}
\newcommand{\notespropositionname}{Proposition}
\newcommand{\notescorollaryname}{Corollary}
\newcommand{\notesdefinitionname}{Definition}
\newcommand{\notesexamplename}{Example}
\newcommand{\notesproblemname}{Problem}
\newcommand{\notesexercisename}{Exercise}
\newcommand{\notesremarkname}{Remark}
\newcommand{\notessolutionname}{Solution}
\fi

\usepackage[a4paper,margin=30mm]{geometry}
\usepackage{microtype}
\usepackage{graphicx}
\usepackage{enumitem}

\newcommand{\exerciseheading}[2]{%
  \par\noindent
  {\large\bfseries #1\par}
  \vspace{0.1em}
  \noindent{\large #2\par}
  \vspace{1.5em}
}

\usepackage{amsmath}
\usepackage{amssymb}
\usepackage{amsthm}
\usepackage{mathtools}
\usepackage{aliascnt}

% Row operations beside the affected row, e.g. \rowop{+2R_1}.
\newcommand{\rowop}[1]{\mathrel{\big|}\,#1}

% Linear systems: the argument gives the number of alignment pairs.
\newenvironment{linearsystem}[1]
{\left\{
  \begin{alignedat}{#1}}
    {
  \end{alignedat}\right.}

  \usepackage[hidelinks,pdfusetitle]{hyperref}

  \numberwithin{equation}{section}
  \setlist{itemsep=0.25em,topsep=0.5em}

  \theoremstyle{plain}
  \newtheorem{theorem}{\notestheoremname}[section]
  \newaliascnt{lemma}{theorem}
  \newtheorem{lemma}[lemma]{\noteslemmaname}
  \aliascntresetthe{lemma}
  \newaliascnt{proposition}{theorem}
  \newtheorem{proposition}[proposition]{\notespropositionname}
  \aliascntresetthe{proposition}
  \newaliascnt{corollary}{theorem}
  \newtheorem{corollary}[corollary]{\notescorollaryname}
  \aliascntresetthe{corollary}

  \theoremstyle{definition}
  \newaliascnt{definition}{theorem}
  \newtheorem{definition}[definition]{\notesdefinitionname}
  \aliascntresetthe{definition}
  \newaliascnt{example}{theorem}
  \newtheorem{example}[example]{\notesexamplename}
  \aliascntresetthe{example}
  \newaliascnt{problem}{theorem}
  \newtheorem{problem}[problem]{\notesproblemname}
  \aliascntresetthe{problem}
  \newtheorem{exercise}{\notesexercisename}

  \theoremstyle{remark}
  \newaliascnt{remark}{theorem}
  \newtheorem{remark}[remark]{\notesremarkname}
  \aliascntresetthe{remark}

  \newenvironment{solution}
  {
  \begin{proof}[\notessolutionname]}
    {
  \end{proof}}

  \usepackage[nameinlink,noabbrev]{cleveref}

  \crefname{theorem}{theorem}{theorems}
  \crefname{lemma}{lemma}{lemmas}
  \crefname{proposition}{proposition}{propositions}
  \crefname{corollary}{corollary}{corollaries}
  \crefname{definition}{definition}{definitions}
  \crefname{example}{example}{examples}
  \crefname{problem}{problem}{problems}
  \crefname{exercise}{exercise}{exercises}
  \crefname{remark}{remark}{remarks}

  \ifx\notelanguage\notesfinnishlanguage
  \crefname{theorem}{lause}{lauseet}
  \crefname{lemma}{lemma}{lemmat}
  \crefname{proposition}{propositio}{propositiot}
  \crefname{corollary}{korollaari}{korollaarit}
  \crefname{definition}{määritelmä}{määritelmät}
  \crefname{example}{esimerkki}{esimerkit}
  \crefname{problem}{tehtävä}{tehtävät}
  \crefname{exercise}{tehtävä}{tehtävät}
  \crefname{remark}{huomautus}{huomautukset}
  \fi

  \newcommand{\NN}{\mathbb{N}}
  \newcommand{\NNone}{\mathbb{N}_1}
  \newcommand{\ZZ}{\mathbb{Z}}
  \newcommand{\QQ}{\mathbb{Q}}
  \newcommand{\RR}{\mathbb{R}}
  \newcommand{\CC}{\mathbb{C}}
  \newcommand{\PP}{\mathbb{P}}

  \DeclareMathOperator{\im}{im}
  \DeclareMathOperator{\rank}{rank}
  \DeclareMathOperator{\syt}{syt}
  \DeclareMathOperator{\pyj}{pyj}
  \DeclareMathOperator{\dom}{dom}
  \DeclarePairedDelimiter{\abs}{\lvert}{\rvert}
  \DeclarePairedDelimiter{\norm}{\lVert}{\rVert}
  \DeclarePairedDelimiterXPP{\maxnorm}[1]{}{\lVert}{\rVert}{_\infty}{#1}
  \newcommand{\set}[1]{\left\lbrace#1\right\rbrace}
  \newcommand{\maxset}[1]{\max\set{#1}}
  \newcommand{\minset}[1]{\min\set{#1}}
  \newcommand{\powerset}[1]{\mathcal{P}\!\left(#1\right)}


\theoremstyle{definition}
\newtheorem{groupexercise}{Ryhmätyö}

\title{Analyysi II\\Harjoitus 5}
\author{}
\date{}

\begin{document}

\exerciseheading{Analyysi II}{Harjoitus 5}

\begin{exercise}
  Osoita, että \(\diffop{x} \cos x = -\sin x\) kaikilla \(x \in \RR\).
\end{exercise}

\begin{solution}
  Käytetään tunnettua raja-arvoa \(\lim_{h \to 0} \frac{\sin h}{h} = 1\).
  Kosinin jatkuvuuden nojalla \(\lim_{h \to 0}(1 + \cos h) = 2\).
  Nimittäjä \(1 + \cos h\) on siis nollasta poikkeava nollan lähellä, joten
  \[
    \lim_{h \to 0} \frac{1 - \cos h}{h}
    = \lim_{h \to 0} \frac{1 - \cos^2 h}{h(1 + \cos h)}
    = \lim_{h \to 0} \frac{h}{1 + \cos h}\left(\frac{\sin h}{h}\right)^2
    = 0.
  \]

  Olkoon \(x \in \RR\) mielivaltainen. Kosinin yhteenlaskukaavan nojalla
  saadaan, kun \(h \neq 0\),
  \[
    \frac{\cos(x + h) - \cos x}{h}
    = -\cos x\,\frac{1 - \cos h}{h} - \sin x\,\frac{\sin h}{h}.
  \]
  Raja-arvojen laskusääntöjen nojalla
  \[
    \lim_{h \to 0} \frac{\cos(x + h) - \cos x}{h}
    = -\cos x \cdot 0 - \sin x \cdot 1
    = -\sin x.
  \]
  Derivaatan määritelmän nojalla siis \(\diffop{x} \cos x = -\sin x\)
  kaikilla \(x \in \RR\).
\end{solution}

\begin{exercise}
  Laske \(\diffop{x} \tan x\) kaikilla
  \(x \in \RR \setminus \bigl\{\frac{\pi}{2} + n\pi : n \in \ZZ\bigr\}\)
  eli pisteissä, joissa \(\tan x = \frac{\sin x}{\cos x}\) on määritelty.
\end{exercise}

\begin{solution}
  Olkoon \(x\) tangentin määrittelyjoukon mielivaltainen piste.
  Koska sini ja kosini ovat derivoituvia ja \(\cos x \neq 0\),
  saadaan osamäärän derivointisäännöllä
  \[
    \begin{aligned}
      \diffop{x} \tan x &= \diffop{x}\left(\frac{\sin x}{\cos x}\right) \\
      &= \frac{\cos x \cdot \diffop{x} \sin x - \sin x \cdot
      \diffop{x} \cos x}{\cos^2 x} \\
      &= \frac{\cos^2 x + \sin^2 x}{\cos^2 x} \\
      &= \frac{1}{\cos^2 x}.
    \end{aligned}
  \]
  Koska \(x\) oli mielivaltainen, kaava pätee koko tangentin määrittelyjoukossa.
\end{solution}

\begin{exercise}
  Olkoon \(f \colon \RR \to \RR\) derivoituva ja \(n \in \NNone\).
  Osoita, että
  \[
    \diffop{x}\bigl(f^n(x)\bigr) = n f^{n-1}(x) \cdot f'(x)
  \]
  kaikilla \(x \in \RR\).
\end{exercise}

\begin{solution}
  Kun \(n = 1\), määritellään \(f^0(x) = 1\) kaikilla \(x \in \RR\),
  jolloin väite palautuu yhtälöön \(\diffop{x} f(x) = f'(x)\).
  Olkoon siis \(n \geq 2\) ja määritellään
  \(g \colon \RR \to \RR\), \(g(t) = t^n\). Potenssisäännön mukaan
  \(g'(t) = n t^{n-1}\) kaikilla \(t \in \RR\).

  Olkoon \(x \in \RR\) mielivaltainen. Koska \(f^n(x) = (g \circ f)(x)\)
  ja molemmat funktiot ovat derivoituvia, saadaan ketjusäännöllä
  \[
    \begin{aligned}
      \diffop{x}\bigl(f^n(x)\bigr) &= (g \circ f)'(x) \\
      &= g'(f(x)) \cdot f'(x) \\
      &= n f^{n-1}(x) \cdot f'(x).
    \end{aligned}
  \]
  Koska \(x\) oli mielivaltainen, kaava pätee kaikilla \(x \in \RR\).
\end{solution}

\newpage

\begin{exercise}
  Jatkoa harjoituksen~2 tehtävään~3. Tarkastellaan funktion
  \(f(x) = (x^3 - 1)^2\) rajoittumaa välille \((-\infty, 1]\),
  jonka käänteisfunktioksi johdettiin
  \[
    g_1 \colon [0, \infty) \to (-\infty, 1],
    \qquad g_1(y) = \sqrt[3]{1 - \sqrt{y}}.
  \]
  Laske derivaatta \(g_1'(4)\) kahdella eri tavalla: Ensin käyttämällä
  käänteisfunktion derivointikaavaa (Lause~5.2.13) ja sitten
  derivoimalla käänteisfunktio \(g_1\).
\end{exercise}

\begin{solution}
  \emph{Käänteisfunktion derivointikaava.}
  Polynomina \(f\) on derivoituva, ja
  \[
    f'(x) = 6x^2(x^3 - 1),
    \qquad f'(-1) = -12 \neq 0.
  \]
  Käänteisfunktio \(g_1\) on määritelty pisteen \(4\) ympäristössä
  ja on siinä jatkuva juurifunktioiden jatkuvuuden nojalla.
  Sovelletaan lausetta~5.2.13 rajoittumaan \(f|_{(-\infty, 1)}\),
  jonka käänteisfunktio on \(g_1|_{(0, \infty)}\).
  Koska \(g_1(4) = -1\), saadaan
  \[
    g_1'(4) = \frac{1}{f'(g_1(4))} = \frac{1}{f'(-1)} = -\frac{1}{12}.
  \]

  \emph{Suora derivointi.}
  Kun \(y > 1\), voidaan kirjoittaa \(g_1(y) = -(\sqrt{y} - 1)^{1/3}\).
  Ketjusäännöllä saadaan
  \[
    \begin{aligned}
      g_1'(y) &= -\frac{1}{3}(\sqrt{y} - 1)^{-2/3} \cdot \frac{1}{2\sqrt{y}} \\
      &= -\frac{1}{6\sqrt{y}(\sqrt{y} - 1)^{2/3}}.
    \end{aligned}
  \]
  Sijoittamalla \(y = 4\) saadaan
  \[
    g_1'(4) = -\frac{1}{6 \cdot 2 \cdot 1^{2/3}} = -\frac{1}{12}.
    \qedhere
  \]
\end{solution}

\begin{exercise}
  Olkoon \(f \colon (0, \infty) \to (0, \infty)\), \(f(x) = x^x\).\newline
  Derivoi funktio \(f\).
\end{exercise}

\begin{solution}
  Kun \(x > 0\), voidaan kirjoittaa \(f(x) = x^x = e^{x\ln x}\).
  Ketju- ja tulosäännöllä saadaan
  \begin{align*}
    f'(x) &= e^{x\ln x} \cdot \diffop{x}(x\ln x) \\
    &= e^{x\ln x} \left(\ln x + x \cdot \frac{1}{x}\right) \\
    &= x^x(\ln x + 1). \qedhere
  \end{align*}
\end{solution}

\begin{exercise}
  Osoita, että funktio \(f \colon \RR \to \RR\),
  \[
    f(x) =
    \begin{cases}
      x\sin\frac{1}{x}, & \text{jos } x \neq 0, \\
      0, & \text{jos } x = 0,
    \end{cases}
  \]
  on origossa jatkuva mutta ei derivoituva.
\end{exercise}

\begin{solution}
  Olkoon \(\varepsilon > 0\) ja valitaan \(\delta = \varepsilon\).
  Kun \(0 < \abs{x} < \delta\), saadaan \(\abs{f(x) - f(0)} \leq \abs{x}
  < \varepsilon\), koska \(\abs{\sin(1/x)} \leq 1\).
  Kun \(x = 0\), erotus on nolla. Siis \(f\) on jatkuva origossa.

  Olkoon \(g \colon \RR \setminus \set{0} \to \RR\), \(g(x) = \sin(1/x)\).
  Kun \(x \neq 0\), erotusosamäärä origossa on \(\frac{f(x)-f(0)}{x} = g(x)\).
  Määritellään lukujonot \((x_n)\) ja \((y_n)\) asettamalla kaikilla
  \(n \in \NNone\)
  \[
    x_n = \frac{1}{2\pi n + \pi/2}, \qquad
    y_n = \frac{1}{2\pi n + 3\pi/2}.
  \]
  Tällöin \(\lim_{n \to \infty} x_n = \lim_{n \to \infty} y_n = 0\).
  Sinin jaksollisuuden nojalla \(g(x_n) = 1\) ja \(g(y_n) = -1\) kaikilla
  \(n \in \NNone\). Raja-arvon jonokriteerin nojalla \(\lim_{x \to 0} g(x)\)
  ei ole olemassa. Siis erotusosamäärällä ei ole raja-arvoa ja \(f\) ei ole
  derivoituva origossa.
\end{solution}

\end{document}
